\apendice{Plan de Proyecto Software}

\section{Introducción}

En este apéndice se documenta la planificación temporal del proyecto, la metodología de gestión utilizada y el estudio de viabilidad económica y legal que sustenta la realización del trabajo.

El proyecto es un \textbf{Clasificador de señales EEG para la detección de TDAH mediante Machine Learning y Deep Learning}. Su objetivo es construir un sistema de apoyo para la clasificación binaria de señales EEG (ADHD vs Control) mediante una aplicación web que integra modelos de Machine Learning clásico y Deep Learning, con énfasis en el rigor metodológico de la evaluación \emph{cross-subject}.

El producto entregado contiene:

\begin{itemize}
    \item Una API REST con FastAPI.
    \item Una SPA desarrollada con React + Vite.
    \item Una base de datos PostgreSQL para almacenar los entrenamientos.
    \item Un \textbf{pipeline de investigación reproducible} en \ruta{scripts/} para entrenar y exportar los modelos que se sirven a través de la API.
    \item Tests automatizados, CI con GitHub Actions y análisis con SonarCloud.
\end{itemize}

\section{Planificación temporal}

\subsection{Metodología}

Se ha utilizado un enfoque \textbf{Scrum aplicado a un TFG individual}, soportado por \textbf{ZenHub} como gestor de tareas integrado con GitHub Issues. Las características más relevantes del flujo de trabajo aplicado son:

\begin{itemize}
    \item \textbf{Fases temáticas de duración variable} (entre 1 y 3 semanas), ajustadas al ritmo académico y a la disponibilidad real del autor.
    \item \textbf{Backlog estructurado en \emph{issues}} alojadas en GitHub y organizadas en el tablero ZenHub. Las issues se crearon de forma distribuida a lo largo del proyecto (3 en febrero, 12 en marzo, 5 en abril y 14 en mayo de 2026), evidenciando un trabajo iterativo y sostenido.
    \item \textbf{Estimación con Story Points} en ZenHub para cuantificar el esfuerzo relativo del trabajo realizado.
    \item \textbf{Trazabilidad commit $\leftrightarrow$ issue}: los commits relevantes enlazan con las \emph{issues} a las que pertenecen.
    \item \textbf{Revisiones con el tutor} en las reuniones programadas a lo largo del curso.
\end{itemize}

El tablero ZenHub puede consultarse en \url{https://app.zenhub.com/workspaces/tfg-adhd-class-69bc5516af887f0036c68928/board}. Al requerir autenticación, la captura del tablero y el burndown del sprint de cierre se incluyen como figuras de este apéndice. Adicionalmente se invitará al usuario \texttt{ubutfgm} y a los miembros del tribunal para acceso directo.

\subsection{Cronograma y fases del desarrollo}

El proyecto se ha desarrollado de forma continua entre el \textbf{26 de febrero de 2026} y la \textbf{fecha prevista de defensa: 6 de julio de 2026}, lo que arroja una duración total de aproximadamente \textbf{19 semanas} (4,5 meses).

El trabajo se ha organizado en ocho fases temáticas seguidas de una fase final dedicada a documentación y entrega. Cada fase corresponde a un incremento funcional verificable en el repositorio mediante \emph{commits} firmados.

\tablaSmallSinColores{Fases del desarrollo del proyecto.}{p{1cm}p{3cm}p{4cm}p{6cm}}{fases-resumen}{
    \textbf{Fase} & \textbf{Fechas} & \textbf{Tema} & \textbf{Resultado verificable} \\
}{
    0 & 26 feb -- 9 mar & Arranque del proyecto & Repositorio creado, dataset descargado, artículos enlazados \\
    1 & 18 mar -- 27 mar & Pipeline base de datos & Carga, preprocesado, epochs, features básicas, CV inicial \\
    2 & 6 abr -- 15 abr & Features avanzadas & Features espectrales, combinadas, CV solo en \emph{train} \\
    3 & 19 abr -- 21 abr & Modelos avanzados & XGBoost añadido; modelos DL (CNN, CNN-LSTM, EEGNet) \\
    4 & 30 abr -- 3 may & Aplicación básica & Exportación del mejor modelo, Streamlit inicial, pestañas \\
    5 & 10 may -- 18 may & Migración a React + FastAPI & SPA React + Vite, Docker Compose, CI, tests, SonarCloud \\
    6 & 18 may -- 27 may & Calidad y robustez & Refactor central, robustez de la API, LICENSE, BD inicial \\
    7 & 27 may -- 30 may & Histórico de experimentos & Esquema BD (Alembic), pestaña ``Experimentos'' \\
    8 & 1 jun -- 6 jul & Memoria, anexos y entrega & Documentación completa, vídeos, release final, defensa \\
}

\subsection{Hitos principales}

\tablaSmallSinColores{Hitos principales del proyecto y su evidencia en el repositorio.}{p{6cm}p{2.5cm}p{4.5cm}}{hitos}{
    \textbf{Hito} & \textbf{Fecha} & \textbf{Evidencia} \\
}{
    Inicio del proyecto (commit inicial) & 26 feb 2026 & Commit \texttt{5064679} \\
    Primer pipeline ML completo con CV & 27 mar 2026 & Commit \texttt{de3a372} \\
    Features espectrales integradas & 6 abr 2026 & Commit \texttt{a587b58} \\
    Modelos DL operativos & 21 abr 2026 & Commit \texttt{e2011f0} \\
    Primera versión de la app (Streamlit) & 30 abr 2026 & Commit \texttt{5ddfc79} \\
    Migración a React + FastAPI completada & 14 may 2026 & Commit \texttt{924d1cb} \\
    CI verde y tests automatizados & 18 may 2026 & Commit \texttt{f1eb637} \\
    Licencia MIT añadida & 25 may 2026 & Commit \texttt{759f1f5} \\
    BD de experimentos en producción & 27 may 2026 & Commit \texttt{c9ca85a} \\
    Pestaña ``Experimentos'' en frontend & 30 may 2026 & Commit \texttt{0afa2bd} \\
    Entrega de memoria y anexos & 6 jul 2026 & Por confirmar \\
}

\subsection{Sprint de consolidación y burndown}

Aunque las fases descritas reflejan la estructura temática del desarrollo, ZenHub permite además organizar las \emph{issues} en sprints \emph{time-boxed} para cuantificar el esfuerzo y obtener métricas Scrum estándar. El último sprint del proyecto (\textbf{May 19 -- Jun 2, 2026}) actúa como \textbf{sprint de consolidación}: agrupa las \emph{issues} completadas a lo largo del proyecto, cada una de ellas estimada previamente con Story Points (SP).

\tablaSmallSinColores{Métricas del sprint de consolidación.}{p{8cm}p{4cm}}{sprint-consolidacion}{
    \textbf{Métrica} & \textbf{Valor} \\
}{
    Rango temporal & May 19 -- Jun 2, 2026 \\
    \emph{Issues} comprometidas & 30 \\
    Story Points comprometidos & 223 \\
    \emph{Issues} completadas & 30 (100\%) \\
    Story Points completados & 223 (100\%) \\
    Cierre efectivo (gráfica) & 27 may 2026 \\
}

El burndown chart generado automáticamente por ZenHub muestra la progresión de los Story Points restantes frente a la línea ideal. La consolidación efectiva del trabajo se produjo cinco días antes del cierre teórico del sprint, evidenciando que la estimación realizada fue conservadora y que el ritmo de cierre permitió completar los 223~SP comprometidos.

% Cuando tengas la captura del burndown, descomenta la siguiente línea y
% asegúrate de colocar el fichero en img/burndown_sprint_may19.png.
%\imagen{burndown_sprint_may19}{Sprint Burndown del sprint de consolidación (May 19 -- Jun 2, 2026). Story Points completados: 223 de 223 (100\%); \emph{Issues} completadas: 30 de 30 (100\%). El cierre efectivo se produjo el 27 de mayo, cinco días antes del fin teórico del sprint.}

% Cuando tengas la captura del tablero, descomenta la siguiente línea y
% asegúrate de colocar el fichero en img/zenhub_board.png.
%\imagen{zenhub_board}{Estado del tablero ZenHub al cierre del proyecto. Se observan los ocho carriles del flujo Kanban y la distribución de las 34 \emph{issues} entre \emph{Review/QA}, \emph{Done} y \emph{Closed}.}

\subsection{Estimación de dedicación}

La dedicación al TFG ha sido \textbf{parcial}, compatibilizándola con otras asignaturas del último curso del Grado. Se estima una dedicación media de \textbf{25--30 horas semanales} durante las 19 semanas activas:

\tablaSmallSinColores{Estimación de dedicación al proyecto.}{p{8cm}p{4cm}}{dedicacion}{
    \textbf{Concepto} & \textbf{Estimación} \\
}{
    Semanas activas & 19 \\
    Horas semanales medias & 27,5 \\
    \textbf{Total estimado} & \textbf{$\approx$ 525 horas} \\
}

Esta estimación incluye el trabajo de investigación, desarrollo de software, redacción de documentación y reuniones con el tutor.

\section{Estudio de viabilidad}

\subsection{Viabilidad económica}

A continuación se desglosan los costes asociados al desarrollo del proyecto. Conforme a las recomendaciones de la guía oficial de TFG de la UBU, se imputan \textbf{costes de personal con la parte impositiva correspondiente a la empresa} y se aplica \textbf{amortización} sobre el hardware (no se imputa el coste total de adquisición).

\subsubsection{Coste de personal}

Se considera el perfil de un \textbf{Ingeniero Informático Junior} desarrollando las labores típicas de un TFG (investigación, desarrollo, documentación, reuniones de seguimiento).

\tablaSmallSinColores{Coste mensual del perfil de Ingeniero Informático Junior.}{p{8cm}p{4cm}}{coste-personal}{
    \textbf{Concepto} & \textbf{Valor} \\
}{
    Salario bruto anual de referencia (junior) & 27.000~€ \\
    Salario mensual bruto (14 pagas) & 1.928,57~€ \\
    Coste empresa adicional ($\sim$32\%) & 617,14~€ \\
    \textbf{Coste mensual total para la empresa} & \textbf{2.545,71~€} \\
    Horas laborables por mes (jornada completa) & 160 \\
    \textbf{Coste por hora (jornada completa)} & \textbf{$\approx$ 15,90~€/h} \\
}

La parte impositiva a cargo de la empresa incluye contingencias comunes (23,60\%), desempleo (5,50\%, tipo general), FOGASA (0,20\%), Formación Profesional (0,60\%), Mecanismo de Equidad Intergeneracional (0,67\%) y AT/EP (variable, estimación $\sim$1,5\%). El total aproximado es $\sim$32\% sobre el salario bruto, conforme a las tablas publicadas por la Seguridad Social para 2025-2026.

\tablaSmallSinColores{Coste de personal imputado al proyecto.}{p{6cm}p{3cm}p{3cm}}{coste-personal-imputado}{
    \textbf{Concepto} & \textbf{Cálculo} & \textbf{Importe} \\
}{
    Horas dedicadas & 525~h & --- \\
    Coste por hora & 15,90~€ & --- \\
    \textbf{Coste total personal} & 525 $\times$ 15,90~€ & \textbf{$\approx$ 8.348~€} \\
}

\subsubsection{Coste de hardware (amortización)}

Se aplica amortización lineal del hardware durante su vida útil, imputando al proyecto únicamente la parte proporcional al tiempo de uso (4,5 meses).

\tablaSmallSinColores{Coste de hardware imputado con amortización lineal.}{p{4cm}p{1.5cm}p{1.5cm}p{2.5cm}p{1.5cm}p{2cm}}{coste-hw}{
    \textbf{Equipo} & \textbf{Precio} & \textbf{Vida útil} & \textbf{Amort. mensual} & \textbf{Meses} & \textbf{Imputado} \\
}{
    Portátil de desarrollo & 1.500~€ & 5 años & 25,00~€ & 4,5 & 112,50~€ \\
    Monitor externo & 220~€ & 5 años & 3,67~€ & 4,5 & 16,50~€ \\
    \textbf{Total hardware} & --- & --- & --- & --- & \textbf{129~€} \\
}

\subsubsection{Coste de software}

Todas las herramientas y bibliotecas utilizadas son \textbf{software libre} (licencias MIT, BSD, Apache 2.0 y LGPL). No se ha adquirido ninguna licencia comercial.

\tablaSmallSinColores{Coste de software del proyecto.}{p{8cm}p{3cm}}{coste-sw}{
    \textbf{Concepto} & \textbf{Coste} \\
}{
    Sistema operativo (Windows preinstalado en el equipo) & 0~€ \\
    IDE (Visual Studio Code) & 0~€ \\
    Bibliotecas de Python y npm & 0~€ \\
    GitHub (cuenta gratuita) & 0~€ \\
    SonarCloud (plan gratuito open source) & 0~€ \\
    \textbf{Total software} & \textbf{0~€} \\
}

\subsubsection{Coste de infraestructura}

El proyecto se ha desarrollado y probado en local, sin coste de hosting o servicios en la nube de pago.

\tablaSmallSinColores{Coste de infraestructura del proyecto.}{p{8cm}p{3cm}}{coste-infra}{
    \textbf{Concepto} & \textbf{Coste} \\
}{
    Hosting cloud & 0~€ \\
    Dominio & 0~€ \\
    \textbf{Total infraestructura} & \textbf{0~€} \\
}

\subsubsection{Resumen de costes}

\tablaSmallSinColores{Resumen total de costes del proyecto.}{p{8cm}p{4cm}}{resumen-costes}{
    \textbf{Categoría} & \textbf{Importe} \\
}{
    Personal & 8.348~€ \\
    Hardware & 129~€ \\
    Software & 0~€ \\
    Infraestructura & 0~€ \\
    \textbf{TOTAL} & \textbf{$\approx$ 8.477~€} \\
}

Los importes son una estimación académica orientativa elaborada conforme a la guía oficial. Las cifras del coste de personal usan valores típicos del mercado español 2025-2026 para un perfil junior y los tipos impositivos vigentes publicados por la Seguridad Social.

\subsection{Viabilidad legal}

\subsubsection{Licencia del software desarrollado}

El proyecto se distribuye bajo \textbf{MIT License} (ver fichero \ruta{LICENSE} en la raíz del repositorio). MIT es una licencia \emph{permisiva de software}, no de contenido, conforme al requisito de la guía oficial UBU, que excluye explícitamente las licencias Creative Commons para el código.

La decisión de licencia se justifica por:

\begin{enumerate}
    \item Es una licencia ampliamente reconocida en la comunidad open source.
    \item Facilita la reutilización académica y profesional del código.
    \item Es compatible con todas las dependencias seleccionadas.
\end{enumerate}

\subsubsection{Análisis de licencias de dependencias}

Se ha realizado un análisis sistemático de la licencia de cada una de las \textbf{30 dependencias directas} del proyecto. Resumen ejecutivo:

\tablaSmallSinColores{Resumen del análisis de licencias por familia.}{p{5cm}p{4cm}p{4cm}}{licencias-familia}{
    \textbf{Familia de licencia} & \textbf{Nº de dependencias} & \textbf{Compat. con MIT} \\
}{
    MIT & 19 & Total \\
    BSD-3-Clause & 7 & Total \\
    Apache 2.0 & 3 & Total \\
    LGPL-3.0-or-later & 1 (\texttt{psycopg}) & Total (uso como librería) \\
    \textbf{TOTAL} & \textbf{30} & \textbf{100\% compatibles} \\
}

Ninguna dependencia es \emph{copyleft} fuerte (GPL o AGPL), por lo que el código propio del TFG puede mantenerse bajo MIT sin restricciones adicionales.

\subsubsection{Tratamiento de datos personales}

El dataset utilizado (EEG ADHD de la Universidad Shahed) es de \textbf{acceso público en IEEE DataPort} y se distribuye con fines de investigación. Los sujetos son anónimos: cada uno se identifica únicamente con un identificador numérico (\texttt{ID}) sin información personal asociada.

La aplicación desarrollada:

\begin{itemize}
    \item \textbf{No persiste datos clínicos crudos} del paciente. Los CSV cargados por el usuario se procesan en memoria y solo se almacenan en la base de datos los metadatos estadísticos del dataset (\emph{hash}, número de filas y columnas, número de sujetos, distribución de clases, canales detectados).
    \item \textbf{No identifica al usuario} de la aplicación (sistema sin autenticación, uso académico).
\end{itemize}

Por tanto, no se procesan datos personales en el sentido del \textbf{Reglamento General de Protección de Datos (RGPD, UE 2016/679)} ni de la \textbf{LOPDGDD 3/2018}. No se requiere una evaluación de impacto adicional.

\subsubsection{Carácter no clínico}

La aplicación tiene \textbf{carácter académico y demostrativo}. No constituye un producto sanitario en el sentido del Reglamento (UE) 2017/745 y no debe utilizarse como herramienta de diagnóstico clínico. Esta limitación se hace explícita en el README y en la memoria del proyecto.
